\section{แบบฝึกหัดเพิ่มเติมสำหรับเตรียมสอบปลายภาค}
\subsection{บทที่ 6 การแปลงลาปลาซ (Laplace Transform)}

\textbf{รายวิชา:} 5571102 คณิตศาสตร์วิศวกรรมไฟฟ้า\\
\textbf{Electrical Engineering Mathematics}

เอกสารชุดฝึกนี้จัดทำขึ้นโดยอ้างอิงแนวทางจาก \textbf{ตัวอย่างที่ 6.3, 6.4, 6.5, 6.6, 6.13} และ \textbf{แบบฝึกหัดท้ายบทข้อ 1, 2, 3, 6, 7} โดยกำหนดให้ระดับความยากของโจทย์ไม่เกินกรอบของตัวอย่างและแบบฝึกหัดต้นฉบับ ชุดฝึกที่ 3 เพิ่มจำนวนขั้นตอนในการจัดรูปและวิเคราะห์ แต่ยังใช้เทคนิคเดิมที่ปรากฏในบทเรียน

\par\noindent\rule{\textwidth}{0.4pt}\par

\section{ชุดฝึกที่ 1}

\subsection{ข้อ 1 — Laplace Transform โดยใช้สมบัติเชิงเส้นและตาราง}

จงหา Laplace transform ของ

\[
f(t)=3+2t+4t^2
\]

\textbf{แนวอ้างอิง:} ตัวอย่าง 6.3 และแบบฝึกหัดท้ายบทข้อ 1

\par\noindent\rule{\textwidth}{0.4pt}\par

\subsection{ข้อ 2 — Inverse Laplace จากตารางโดยตรง}

จงหา inverse Laplace ของ

\[
F(s)=\frac{5}{s}+\frac{3}{s+2}-\frac{2}{s^2}
\]

\textbf{แนวอ้างอิง:} แบบฝึกหัดท้ายบทข้อ 3

\par\noindent\rule{\textwidth}{0.4pt}\par

\subsection{ข้อ 3 — Inverse Laplace โดยใช้ Partial Fractions: รากจริงต่างกัน}

จงหา

\[
\mathcal{L}^{-1}
\left\{
\frac{4s+7}{(s+1)(s+3)}
\right\}
\]

โดยแสดงขั้นตอนการแยกเศษส่วนย่อย

\textbf{แนวอ้างอิง:} ตัวอย่าง 6.5, ตัวอย่าง 6.13 และแบบฝึกหัดท้ายบทข้อ 6

\par\noindent\rule{\textwidth}{0.4pt}\par

\subsection{ข้อ 4 — Inverse Laplace โดยใช้ Partial Fractions: รากซ้ำ}

จงหา

\[
\mathcal{L}^{-1}
\left\{
\frac{s+5}{s(s+2)^2}
\right\}
\]

โดยตั้งรูปเศษส่วนย่อยให้ครบทุกพจน์

\textbf{แนวอ้างอิง:} ตัวอย่าง 6.6 และแบบฝึกหัดท้ายบทข้อ 7

\par\noindent\rule{\textwidth}{0.4pt}\par

\subsection{ข้อ 5 — Inverse Laplace โดยการทำกำลังสองสมบูรณ์}

จงหา

\[
\mathcal{L}^{-1}
\left\{
\frac{2s+7}{s^2+4s+13}
\right\}
\]

โดยจัดตัวส่วนให้อยู่ในรูปกำลังสองสมบูรณ์ก่อนเทียบตาราง

\textbf{แนวอ้างอิง:} ตัวอย่าง 6.4

\par\noindent\rule{\textwidth}{0.4pt}\par

\newpage

\section{ชุดฝึกที่ 2}

\subsection{ข้อ 1 — Laplace Transform ของ Exponential, Sine และ Cosine}

จงหา Laplace transform ของ

\[
f(t)=4e^{-2t}-3\sin 4t+2\cos 4t
\]

\textbf{แนวอ้างอิง:} ตัวอย่าง 6.3 และแบบฝึกหัดท้ายบทข้อ 2

\par\noindent\rule{\textwidth}{0.4pt}\par

\subsection{ข้อ 2 — Inverse Laplace จากตารางโดยตรง}

จงหา inverse Laplace ของ

\[
F(s)=\frac{4}{s}-\frac{5}{s+3}+\frac{6}{s^2}
\]

\textbf{แนวอ้างอิง:} แบบฝึกหัดท้ายบทข้อ 3

\par\noindent\rule{\textwidth}{0.4pt}\par

\subsection{ข้อ 3 — Inverse Laplace โดยใช้ Partial Fractions: รากจริงต่างกัน}

จงหา

\[
\mathcal{L}^{-1}
\left\{
\frac{5s+14}{(s+2)(s+4)}
\right\}
\]

โดยแสดงขั้นตอนการหาสัมประสิทธิ์ของเศษส่วนย่อย

\textbf{แนวอ้างอิง:} ตัวอย่าง 6.5, ตัวอย่าง 6.13 และแบบฝึกหัดท้ายบทข้อ 6

\par\noindent\rule{\textwidth}{0.4pt}\par

\subsection{ข้อ 4 — Inverse Laplace โดยใช้ Partial Fractions: รากซ้ำ}

จงหา

\[
\mathcal{L}^{-1}
\left\{
\frac{3s+4}{s(s+2)^2}
\right\}
\]

โดยตั้งเศษส่วนย่อยในรูปที่ถูกต้องก่อนหาค่าสัมประสิทธิ์

\textbf{แนวอ้างอิง:} ตัวอย่าง 6.6 และแบบฝึกหัดท้ายบทข้อ 7

\par\noindent\rule{\textwidth}{0.4pt}\par

\subsection{ข้อ 5 — Inverse Laplace โดยการทำกำลังสองสมบูรณ์}

จงหา

\[
\mathcal{L}^{-1}
\left\{
\frac{3s+11}{s^2+6s+13}
\right\}
\]

โดยจัดตัวส่วนและตัวเศษให้สามารถเทียบกับรูป sine และ cosine ในตาราง Laplace ได้

\textbf{แนวอ้างอิง:} ตัวอย่าง 6.4

\par\noindent\rule{\textwidth}{0.4pt}\par

\newpage

\section{ชุดฝึกที่ 3 — ระดับท้าทาย}

\subsection{ข้อ 1 — Laplace Transform โดยใช้สมบัติเชิงเส้นและตาราง}

จงหา Laplace transform ของ

\[
f(t)=3t^2-4e^{-2t}+5\sin 3t-2\cos 3t
\]

โดยแสดงการใช้สมบัติเชิงเส้นและคู่การแปลงของแต่ละพจน์

\textbf{แนวอ้างอิง:} ตัวอย่าง 6.3 และแบบฝึกหัดท้ายบทข้อ 1–2

\textbf{จุดที่ยากขึ้น:} มีฟังก์ชันหลายชนิดอยู่ในข้อเดียว แต่ยังเป็นการเทียบตารางโดยตรง

\par\noindent\rule{\textwidth}{0.4pt}\par

\subsection{ข้อ 2 — Inverse Laplace โดยเทียบตารางหลายรูปแบบ}

จงหา

\[
\mathcal{L}^{-1}
\left\{
\frac{6}{s^3}
-\frac{5}{s+4}
+\frac{4s}{s^2+16}
+\frac{12}{s^2+16}
\right\}
\]

โดยแสดงการหา inverse Laplace ของแต่ละพจน์

\textbf{แนวอ้างอิง:} ตัวอย่าง 6.3 และแบบฝึกหัดท้ายบทข้อ 3

\textbf{จุดที่ยากขึ้น:} ต้องจำแนกรูป \(t^n\), exponential, sine และ cosine ภายในข้อเดียว

\par\noindent\rule{\textwidth}{0.4pt}\par

\subsection{ข้อ 3 — Partial Fractions กรณีรากจริงต่างกัน 3 ราก}

จงหา

\[
\mathcal{L}^{-1}
\left\{
\frac{3s^2+3s+6}
{s(s+1)(s+3)}
\right\}
\]

โดยดำเนินการดังนี้

\begin{enumerate}
    \item ตั้งรูปเศษส่วนย่อยให้ถูกต้อง
    \item หาค่าสัมประสิทธิ์ \(A,B,C\)
    \item เขียน \(F(s)\) หลังแยกเศษส่วนย่อย
    \item หา \(f(t)\)
\end{enumerate}

\textbf{แนวอ้างอิง:} ตัวอย่าง 6.5, ตัวอย่าง 6.13 และแบบฝึกหัดท้ายบทข้อ 6

\textbf{จุดที่ยากขึ้น:} ตัวเศษเป็นพหุนามกำลังสองและมีรากต่างกัน 3 ราก แต่ใช้หลักการเดิม

\par\noindent\rule{\textwidth}{0.4pt}\par

\subsection{ข้อ 4 — Partial Fractions กรณีมีรากซ้ำ}

จงหา

\[
\mathcal{L}^{-1}
\left\{
\frac{s^2+9s+8}
{s(s+2)^2}
\right\}
\]

โดยเริ่มจากการตั้ง

\[
\frac{s^2+9s+8}{s(s+2)^2}
=
\frac{A}{s}
+
\frac{B}{s+2}
+
\frac{C}{(s+2)^2}
\]

แล้วหาค่า \(A,B,C\) และหา inverse Laplace

\textbf{แนวอ้างอิง:} ตัวอย่าง 6.6 และแบบฝึกหัดท้ายบทข้อ 7

\textbf{จุดที่ยากขึ้น:} ต้องจัดการทั้งรากธรรมดาและรากซ้ำในโจทย์เดียว

\par\noindent\rule{\textwidth}{0.4pt}\par

\subsection{ข้อ 5 — Inverse Laplace โดยทำกำลังสองสมบูรณ์และจัดตัวเศษ}

จงหา

\[
\mathcal{L}^{-1}
\left\{
\frac{4s+18}
{s^2+6s+13}
\right\}
\]

โดยดำเนินการดังนี้

\begin{enumerate}
    \item จัดตัวส่วนให้อยู่ในรูป
    \[
    (s+a)^2+\omega^2
    \]
    \item จัดตัวเศษให้อยู่ในรูปที่สัมพันธ์กับ \(s+a\)
    \item แยกเป็นพจน์ cosine และ sine
    \item หา inverse Laplace
\end{enumerate}

\textbf{แนวอ้างอิง:} ตัวอย่าง 6.4

\textbf{จุดที่ยากขึ้น:} นักศึกษาต้องจัดทั้งตัวส่วนและตัวเศษเองก่อนเทียบตาราง

\par\noindent\rule{\textwidth}{0.4pt}\par


